\documentclass[../../../../SoftwareRequirementsSpecification.tex]{subfiles}
\begin{document}
% Richiede le macro definite in src/macros/useCaseMacros.tex
% (incluse nel preambolo di SoftwareRequirementsSpecification.tex)

\ucstart
\begin{tabularx}{\textwidth}{|l|X|}
  \hline
  \ucid{A-04} \\ \hline
  \ucname{Visualizzazione Schede Assegnate} \\ \hline
  \ucactor{Atleta} \\ \hline
  \ucdescription{L'atleta consulta le schede di
    allenamento che gli sono state assegnate dal proprio PT, fino al
    dettaglio degli esercizi previsti in ciascuna settimana, e li
    confronta con quanto ha registrato di avere svolto.} \\ \hline
  \ucprecond{L'atleta deve essere autenticato.} \\ \hline
  \ucpostcond{L'atleta visualizza gli allenamenti
    previsti per la settimana scelta e, dove presenti, gli
    allenamenti che vi ha registrato.} \\ \hline
  \ucmainflow{
    \item L'atleta atterra sulla pagina delle proprie schede.
    \item Il sistema mostra l'elenco delle schede di allenamento
      assegnate all'atleta, con nome, periodo di validità e stato.
    \item L'atleta preme su una scheda dell'elenco.
    \item Il sistema mostra il dettaglio della scheda assegnata, che
      riporta il nome, il periodo di validità e l'elenco delle
      settimane di allenamento, con evidenziata la settimana corrente
      e segnalate le settimane in cui l'atleta ha registrato almeno
      un allenamento.
    \item L'atleta preme su una settimana, scelta fra tutte quelle
      della scheda, comprese quelle già trascorse (vedi Nota 4).
    \item Il sistema mostra le sessioni di allenamento della
      settimana con gli esercizi previsti e i valori prescritti a
      quell'atleta (vedi Nota 1); per le sessioni in cui l'atleta ha
      registrato un allenamento, accanto ai valori prescritti sono
      riportati quelli che ha effettivamente svolto e gli esercizi
      che ha marcato come saltati (vedi Nota 3).
  } \\ \hline
  \ucaltflow{
    \textbf{2a. Nessuna scheda assegnata} \newline
    - Il sistema mostra un messaggio ("Nessuna scheda di allenamento
    assegnata"). \newline
    - Il flusso termina. \newline
    \textbf{6a. Modifica di un allenamento registrato} \newline
    - L'atleta preme su una sessione per cui ha registrato un
    allenamento. \newline
    - Il flusso continua dal flusso alternativo 5a del caso d'uso
    A-02 (vedi Nota 2).
  } \\ \hline
  \ucnote{
    \textbf{Nota 1:} i valori mostrati al passo 6 non sono
    memorizzati: il sistema li calcola a partire dagli allenamenti
    configurati nel blocco mensile (PT-04), dalla progressione di
    ciascun esercizio, dai pesi inseriti dal PT al momento
    dell'assegnazione (PT-05) e dalle personalizzazioni
    eventualmente introdotte dal PT per quell'atleta
    (PT-09).\newline
    \textbf{Nota 2:} la registrazione di un nuovo allenamento non
    parte da questa pagina, ma dal caso d'uso A-02, che la guida a
    partire dalla data e ne ricava la settimana. Da qui si raggiunge
    soltanto la modifica di un allenamento già registrato, per il
    quale data e sessione sono già note.\newline
    \textbf{Nota 3:} l'assenza di un allenamento registrato significa
    che l'atleta non lo ha registrato, non che non lo ha svolto: il
    sistema non può distinguere i due casi.\newline
    \textbf{Nota 4:} tutto ciò che questo caso d'uso mostra è in sola
    lettura, comprese le settimane trascorse. A differenza del PT
    (caso d'uso PT-09), l'atleta non modifica in alcun caso ciò che
    gli è prescritto.
  } \\ \hline
\end{tabularx}
\end{document}
